\chapter{Una mappa logica della ricerca preistorica in Veneto}
\label{cap:research_bias}

La mappa della preistoria veneta che il capitolo precedente ha costruito provincia per provincia non è una mappa distorta: è una mappa parziale. La differenza è importante. Una mappa distorta riflette la realtà ma con proporzioni alterate; una mappa parziale ha zone illuminate con buona precisione e zone semplicemente buie, in cui il silenzio non dice nulla sulla realtà sottostante. Nelle zone illuminate, dove la ricerca è stata sistematica e continuativa, i risultati sono proporzionali all'investimento e sostanzialmente affidabili. Nelle zone buie, l'assenza di dati non è un dato: è l'assenza di chi ha cercato.

Questo capitolo non parla di distorsioni né di errori della ricerca passata. Parla della logica con cui la ricerca si è distribuita sul territorio veneto nel corso degli ultimi cento anni: perché alcune aree sono state illuminate, chi le ha illuminate, cosa le ha rese candidate alla ricerca sistematica; e perché altre aree sono rimaste nell'ombra, per ragioni che di solito non hanno nulla a che vedere con la loro importanza preistorica e molto a che vedere con la storia delle istituzioni, con la geografia logistica, e con eventi che hanno reindirizzato l'attenzione collettiva.

% ============================================================
\section{Dove la luce è caduta, e perché}
% ============================================================

\subsection{La scuola ferrarese e la Lessinia}

Il fatto più importante per capire la distribuzione della ricerca preistorica nel Veneto non è geografico: è istituzionale. L'Università di Ferrara, attraverso il suo dipartimento di preistoria, è la presenza dominante nella storia della ricerca preistorica veneta per quasi settant'anni. Alberto Broglio ha costruito a partire dagli anni sessanta un programma di ricerca sistematico attorno ai ripari della Lessinia che ha prodotto la sequenza paleolitica più completa dell'Italia nordorientale. Riparo Tagliente fu il sito-guida: una volta documentata la sua straordinaria continuità occupazionale dal Musteriano all'Epigravettiano, attirò finanziamenti, dottorati, collaborazioni internazionali, e spinse verso prospezioni nei siti vicini (Mezzena, Ghiacciaia, San Bernardino). Ogni grande scoperta in un'area aumenta la probabilità che vengano cercate le aree adiacenti: è un meccanismo noto in tutte le scienze empiriche, e in archeologia lo si vede funzionare con precisione nella Lessinia degli anni settanta e ottanta.

Lo stesso programma, esteso e rinnovato da Marco Peresani (allievo di Broglio), si è spostato al Cansiglio nel 1993 e ha prodotto in trent'anni di lavoro sistematico la sequenza epigravettiana-mesolitica più completa delle Prealpi venete: oltre 4.000 manufatti, archivi pollinici, siti di riferimento metodologico per tutta l'archeologia prealpina italiana. Al Cansiglio era già andata l'équipe di Broglio negli anni precedenti per le prime prospezioni. La continuità istituzionale è il filo: stessa scuola, stesso approccio metodologico, risultati comparabili per rigore e densità.

È la stessa scuola ferrarese, con Attilio Guerreschi prima e Federica Fontana poi, che ha scavato Mondeval de Sora a partire dal 1986 e ha scoperto la seconda sepoltura mesolitica al Passo Giau nel 2019. È ancora Broglio che ha diretto gli scavi di Fumane a partire dagli anni novanta, portando alla luce le figure aurignaziane che oggi collocano la grotta tra i siti di riferimento mondiale per l'origine dell'arte figurativa. Il dato ricorrente è questo: dove la scuola ferrarese ha portato il suo programma sistematico, ha trovato risultati di livello internazionale. Dove non è arrivata, la mappa rimane bianca.

\subsection{La Grotta di Fumane: una scoperta che ne genera altre}

Fumane merita un paragrafo separato perché illustra un meccanismo distinto da quello della Lessinia. Non era un'area già nota per i suoi siti paleolitici: la grotta era conosciuta ma non scavata sistematicamente. Fu il programma di Broglio, nel quadro più ampio delle indagini sulla Valpolicella, a portarla alla luce nella sua straordinaria rilevanza. Una volta pubblicati i frammenti dipinti aurignaziani, la grotta è entrata nel circuito internazionale dei siti di riferimento per l'arte paleolitica: le datazioni del 2025 con metodi bayesiani sono il prodotto diretto di quella visibilità, che ha portato collaborazioni con laboratori francesi e spagnoli impossibili senza il profilo internazionale acquisito.

Il meccanismo è circolare ma non vizioso: un sito importante attrae risorse che permettono di documentarlo meglio, che aumentano la sua importanza, che attraggono più risorse. Il problema è che questo circolo virtuoso si innesca solo se qualcuno ha già fatto la prima mossa, cioè ha dedicato tempo e denaro a uno scavo in un'area senza garanzie di risultato. Nel caso di Fumane, quella mossa è stata fatta da un programma universitario già finanziato per altri motivi. Senza quel programma, la grotta sarebbe probabilmente rimasta dove era.

\subsection{Il Monte Luppia: un uomo, una visita, un catalogo}

La storia del Monte Luppia è diversa da tutte le precedenti, e per questo è la più utile per capire le zone buie della mappa. Nel 1964 Mario Pasotti, insegnante, visitò le incisioni della Valcamonica. Al ritorno, riconobbe sulle rocce del Monte Luppia, vicino a casa sua, gli stessi segni che aveva appena visto altrove. Non era un ricercatore professionale, non aveva un programma istituzionale alle spalle, non disponeva di fondi universitari. Aveva uno sguardo addestrato dall'esperienza diretta con un sito già noto.

Fabio Gaggia, che ha poi costruito il catalogo sistematico del complesso e nel 1985 ha fondato un centro di studi dedicato, ha continuato questo lavoro per decenni, producendo il corpus di oltre duecentocinquanta rocce e tremila figure che è oggi il riferimento principale per il sito. Ma il centro di studi di Torri del Benaco non ha mai avuto la struttura, il personale e i finanziamenti di un dipartimento universitario. Il suo raggio di azione è rimasto essenzialmente il territorio benacense. Le Prealpi vicentine, i versanti del Grappa, la Lessinia calcaree al di fuori dei siti già noti di Tagliente e Mezzena: non sono mai stati oggetto di un programma di rilevamento dell'arte rupestre comparabile per sistematicità a quello che, negli stessi decenni, veniva condotto in Valcamonica o in Valle d'Aosta.

La domanda implicita è semplice: se il Monte Luppia è stato scoperto perché qualcuno con lo sguardo giusto ci ha passato, quante superfici analoghe nel Veneto non sono mai state guardate con lo stesso sguardo?

\subsection{Este, Frattesina, Altino: la ricchezza che si trova da sola}

Le scoperte dell'età del ferro e del bronzo nel Veneto meridionale seguono una logica diversa ancora. Este è nota da secoli: la ricchezza dei materiali dalle necropoli (situle, bronzi figurati, corredi funerari elaborati) l'ha resa oggetto di attenzione antiquaria già nell'Ottocento, molto prima che il concetto di scavo sistematico si affermasse. Le prime collezioni sistematiche risalgono agli anni 1870-1880, e il museo di Este è uno dei più antichi del Veneto. Frattesina è stata scoperta in modo più recente ma per ragioni analoghe: una densità di materiali da superficie così eccezionale (avorio, vetro, ambra baltica) da rendere impossibile ignorarla. Altino è emersa progressivamente da decenni di costruzione e bonifica.

Queste scoperte non richiedevano uno sguardo addestrato all'arte rupestre: richiedevano di guardare quello che emergeva dal terreno durante i lavori agricoli, e di riconoscerne il valore. Sono le scoperte che si trovano quasi da sole, perché i materiali sono abbondanti, visibili e inequivocabili. Il contrasto con l'arte rupestre è totale: un'incisione su calcare richiede condizioni di luce specifica, un occhio allenato, spesso strumentazione dedicata. Una situla di bronzo non richiede niente di tutto questo.

% ============================================================
\section{Le zone buie, e perché sono buie}
% ============================================================

\subsection{L'assenza dell'ancora istituzionale}

La variabile più predittiva della densità di siti noti in un'area è l'esistenza di un programma di ricerca continuativo con un'istituzione alle spalle. Dove c'è stata un'università, un museo con un direttore attivo, o anche un singolo ricercatore con la tenacia di Gaggia, la mappa si illumina. Dove non c'è stata nessuna di queste cose, la mappa è bianca.

Il Massiccio del Grappa è il caso più evidente. Morfologicamente è un candidato interessante per la frequentazione preistorica: creste, selle, versanti a quota variabile, vicinanza alla pianura veneta, con le sue necropoli venetiche a pochi chilometri. Ma non è mai stato oggetto di una campagna di prospezione preistorica sistematica comparabile a quelle condotte sulla Lessinia o nel Cadore. Nessuna università ha mai avuto un programma centrato sul Grappa per la preistoria. Nessun Gaggia locale ha mai fondato un centro di studi sulle rocce del Grappa. Il risultato è il silenzio, che non deve essere confuso con l'assenza.

Lo stesso vale per la fascia prealpina trevigiana (Col Visentin, Cesen), per gran parte del Bellunese fuori dal Cadore e dal Cansiglio, per quasi tutto l'arco prealpino vicentino fuori dalla Val d'Assa. Sono aree senza un'ancora istituzionale, dove la ricerca è rimasta episodica, affidata a segnalazioni occasionali o a ricognizioni di superficie mai sistematizzate.

\subsection{Le zone di guerra: un problema doppio}

Il Pasubio, il Carega, Tonezza del Cimone, l'Altopiano di Asiago, la zona del Grappa prealpino, la Marmolada: la Prima Guerra Mondiale ha trasformato fisicamente questi massicci in modo irreversibile, e ha prodotto un effetto secondario altrettanto importante per la ricerca preistorica: ha catturato l'attenzione collettiva e istituzionale per decenni.

Il problema fisico è reale e viene trattato in dettaglio nel capitolo successivo: trincee scavate nelle creste calcaree, strade militari costruite su versanti vergini, minature e crolli controllati, depositi di materiale bellico che rendono tuttora pericolose alcune zone. Tutto questo ha alterato superfici che erano rimaste intatte per millenni.

Ma il problema culturale è forse altrettanto grave. Dopo il 1918, questi massicci sono entrati nell'immaginario collettivo come luoghi della Grande Guerra, non come luoghi della preistoria. I musei dedicati all'Altipiano di Asiago sono musei della Prima Guerra Mondiale. Le guide turistiche del Pasubio descrivono le gallerie e le trincee. I fondi pubblici per la ricerca storica in queste aree sono andati a quella storia, non alla preistoria. L'attenzione istituzionale è stata assorbita dall'evento più recente e più drammatico, e ha lasciato la preistoria in un secondo piano da cui non è mai uscita.

L'unica finestra affidabile sullo stato di queste superfici prima della guerra sono le fonti anteriori al 1915: i bollettini del Club Alpino Italiano, le relazioni delle Società di Scienze Naturali, le guide alpine degli anni 1870-1914. Sono fonti che non cercavano arte rupestre in senso moderno, ma che descrivevano morfologie, ripari, creste, e talvolta segni curiosi su rocce, con la precisione del naturalista d'epoca. Nessuno le ha ancora sistematicamente rilette con gli occhi di chi cerca arte preistorica.

\subsection{L'alta Carnia e le aree di confine}

Sappada, la Val Visdende, l'alta Carnia (Sauris, Ampezzo, Forni) condividono con le zone di guerra un carattere diverso ma ugualmente limitante per la ricerca: sono aree di confine politico, rimaste a lungo in una zona grigia tra competenze istituzionali diverse. La preistoria della Carnia è studiata principalmente nell'ambito della preistoria friulana; quella del Cadore bellunese nell'ambito di quella veneta; quella della Pusteria nell'ambito di quella altoatesina. Le valli che collegano questi sistemi, e che in preistoria erano probabilmente le più frequentate proprio perché connettevano i diversi versanti, ricadono nelle giunture tra le competenze istituzionali e tendono a essere studiate meno di quanto la loro posizione geografica giustificherebbe. Il Lesachtal austriaco, documentato solo in letteratura di lingua tedesca non ancora sistematicamente recepita in italiano, è il caso estremo di questa frammentazione.

\subsection{Il problema specifico dell'arte rupestre}

L'arte rupestre richiede metodi di ricerca diversi da quelli usati per trovare un accampamento o una necropoli, e questo crea una lacuna specifica che si sovrappone a tutte le altre.

Trovare un sito con industria litica richiede di percorrere il terreno guardando il suolo: è un'abilità che ogni archeologo addestrato possiede. Trovare un'incisione su calcare richiede condizioni di luce molto specifiche (luce radente, di preferenza bassa sul versante, mai zenitale), occhio allenato a distinguere il segno intenzionale dall'erosione naturale, e spesso strumentazione dedicata: la riflettografia a trasformata d'immagine (RTI), il rilievo fotogrammetrico tridimensionale, la documentazione con luce strutturata sono oggi le tecniche standard per documentare superfici incise in modo che resista alla revisione critica. Nessuna di queste tecniche è standard nell'arsenale dell'archeologo generalista.

In Italia, i centri con competenza specifica nell'arte rupestre si contano su poche mani: il Centro Camuno di Studi Preistorici a Capo di Ponte, il MUSE di Trento, il gruppo di ricerca coordinato da Angelo Fossati al Politecnico di Milano, poche altre équipe. La Valcamonica ha un'istituzione dedicata che lavora dal 1964. Il Veneto non ha mai avuto niente di equivalente. Gaggia a Torri del Benaco è il riferimento per il Garda, ma il suo lavoro non è mai diventato un programma istituzionale esteso al territorio veneto nella sua interezza.

Il risultato concreto è che non sappiamo se l'arte rupestre sia assente nel Veneto prealpino o semplicemente non sia stata cercata con i metodi giusti. Sono due ipotesi empiricamente indistinguibili nello stato attuale della ricerca, e questo è il silenzio più significativo che la mappa ci consegna.

% ============================================================
\section{La correlazione che conta}
% ============================================================

Mettendo insieme la mappa di dove si è cercato e la mappa di cosa si è trovato, emerge una correlazione che sarebbe sbagliato ignorare: le aree con ricerca sistematica hanno prodotto risultati di rilievo internazionale (Lessinia, Fumane, Cansiglio, Cadore, Monte Luppia, Este, Frattesina); le aree senza ricerca sistematica producono silenzio. La correlazione è quasi perfetta.

Questo non significa che sotto ogni zona buia ci sia necessariamente un sito importante. Alcune zone sono buie per ragioni che hanno anche a che fare con la loro storia preistorica reale: il Polesine non ha arte rupestre non per assenza di ricerca ma per assenza di rocce affioranti; le zone di guerra hanno subito distruzioni fisiche reali; alcune aree alpine possono davvero essere state meno frequentate di altre per ragioni climatiche o ecologiche. La correlazione non è una prova, è un segnale.

Ma è un segnale abbastanza forte da giustificare un'inversione nell'onere della prova. Nel ragionamento corrente, chi afferma che il Veneto potrebbe contenere arte rupestre non ancora trovata deve giustificare questa ipotesi contro l'evidenza del silenzio. Nel ragionamento corretto, dopo aver visto il capitolo precedente e questo, è il contrario: chi conclude che il silenzio equivale all'assenza deve spiegare perché la correlazione tra ricerca e risultati, perfetta ovunque tranne che nelle zone buie, dovrebbe interrompersi proprio là.

\begin{table}[htbp]
\centering
\small
\begin{tabular}{@{}p{3.2cm}p{2.2cm}p{2.8cm}p{2.2cm}@{}}
\toprule
\textbf{Area} & \textbf{Tipo di ricerca} & \textbf{Istituzione di riferimento} & \textbf{Risultato} \\
\midrule
Lessinia & sistematica, continuativa & Università di Ferrara & sequenza di riferimento EU \\
Fumane & sistematica, continuativa & Università di Ferrara & sito mondiale Aurignaziano \\
Cansiglio & sistematica, continuativa & Università di Ferrara & 4.000+ manufatti, 30 anni \\
Cadore/Mondeval & sistematica, poi continuativa & Università di Ferrara & sepolture mesolitiche \\
Monte Luppia & sistematica locale & Centro studi Gaggia & 250 rocce, 3.000 figure \\
Val d'Assa & episodica, incompleta & nessuna & corpus incerto, non finito \\
Este/Frattesina & sistematica & Università di Ferrara & siti di livello europeo \\
Grappa & assente & nessuna & silenzio \\
Prealpi vicentine & assente/episodica & nessuna & silenzio \\
Prealpi trevigiane & assente & nessuna & silenzio \\
Zoldano & assente & nessuna & silenzio (dichiarato) \\
Zone di guerra & assente per preistoria & musei WWI & silenzio \\
\bottomrule
\end{tabular}
\caption{Distribuzione della ricerca preistorica nel Veneto e nei territori contigui, con l'istituzione di riferimento e il risultato prodotto. La correlazione tra tipo di ricerca e risultato è il dato più rilevante della tabella.}
\label{tab:mappa_ricerca}
\end{table}

% ============================================================
\section{Quattro tipi di silenzio}
% ============================================================

Il silenzio della mappa non è uniforme. Vale la pena distinguere almeno quattro tipi diversi, perché hanno implicazioni diverse per la ricerca futura e per le conclusioni che è lecito trarne.

Il \emph{silenzio strutturale} è quello del Polesine o della laguna veneziana: non c'è arte rupestre perché non ci sono rocce affioranti, e non ci sarà mai, qualunque sia l'intensità della ricerca. È il tipo di silenzio più informativo: dice qualcosa di certo, e non deve essere cercato.

Il \emph{silenzio tafonomico} è quello delle creste di guerra, delle cave, delle aree agricolizzate in profondità: poteva esserci qualcosa, ma le condizioni di conservazione erano così sfavorevoli da rendere la sopravvivenza improbabile. Non è certo come il silenzio strutturale, ma nemmeno è colmabile con più ricerca: ciò che è stato distrutto non può essere recuperato. Il capitolo che segue questo dedica un'analisi specifica a questo tipo di silenzio.

Il \emph{silenzio da assenza di ricerca} è quello del Grappa, dello Zoldano, della fascia prealpina trevigiana: nessuno ha cercato sistematicamente, e il silenzio non dice nulla sulla realtà sottostante. È il tipo di silenzio più interessante per la ricerca futura, perché è l'unico in cui nuova ricerca può fare la differenza.

Il \emph{silenzio da metodo inadeguato} è specifico dell'arte rupestre: aree dove si è cercata la frequentazione preistorica in senso generale (con survey di superficie per la litica, con sondaggi per le stratigrafie) ma senza le tecniche specifiche per il rilevamento dell'arte rupestre. In queste aree, il silenzio sull'arte non dice nulla perché è stato usato lo strumento sbagliato per la domanda giusta.

La differenza pratica tra questi quattro tipi è netta. Il silenzio strutturale chiude la domanda. Il silenzio tafonomico la ridimensiona. Il silenzio da assenza di ricerca la mantiene aperta. Il silenzio da metodo inadeguato la rimanda a nuove indagini con strumenti diversi.

% ============================================================
\section{Il confronto che illumina la lacuna}
% ============================================================

Due casi documentati permettono di calibrare meglio quanto pesi la differenza tra ricerca assente e ricerca sistematica, senza limitarsi alla logica interna al Veneto.

Il primo è il progetto Silvretta Historica nelle Alpi svizzere centrali e sudorientali, documentato da Marcel Cornelissen e Thomas Reitmaier \citep{cornelissen2016}. Prima del progetto, le Alpi svizzere centrali erano descritte nella letteratura come un'area con presenza mesolitica rada o assente. In otto anni di ricerca sistematica, il progetto ha censito oltre 230 siti sopra i 2.000 metri in un'area di 550 chilometri quadrati, con una presenza mesolitica che gli autori definiscono ``diffusa e quasi continua''. Non è un ribaltamento della realtà: è la realtà che emerge quando viene cercata con i metodi giusti per il tempo giusto.

Il secondo è lo studio di Davide Visentin e Francesco Carrer sulle Dolomiti venete, pubblicato nel 2017 \citep{visentin2017}, che ha analizzato con modelli statistici di punto quanto della distribuzione nota dei siti mesolitici dipenda dal comportamento insediativo antico e quanto dallo sforzo di ricerca. Il risultato è più sfumato del primo caso: i due fattori, ambientale e da sforzo di ricerca, spiegano in misura pressoché uguale la distribuzione osservata. Né tutto bias, né tutto distribuzione reale: una combinazione dei due, in proporzioni che solo la ricerca sistematica nelle zone ancora bianche potrà chiarire.

Questi due casi non dimostrano che il Veneto nasconda una densità di arte rupestre paragonabile alla Valcamonica. Dimostrano che la correlazione tra intensità della ricerca e risultati è empiricamente documentata, non solo teoricamente plausibile; e che l'inversione del rapporto causa-effetto (l'assenza di risultati come causa, e non come effetto, della scarsità di ricerca) è un errore documentato, non solo un rischio astratto.

% ============================================================
\section{Una distinzione che vale la pena tenere ferma}
% ============================================================

I due casi discussi nell'ultima sezione, come la maggior parte della ricerca citata nel capitolo precedente, documentano la frequentazione preistorica attraverso industria litica, fauna, resti ossei, strutture. Non documentano arte rupestre. Il processo che porta a trovare un accampamento mesolitico è diverso da quello che porta a trovare un pannello inciso: diversi strumenti, diversi criteri di riconoscimento, diverse condizioni ambientali favorevoli.

Questa distinzione non indebolisce il ragionamento di questo capitolo: lo precisa. Il silenzio sull'arte rupestre nel Veneto prealpino non è spiegato solo da ciò che non è stato trovato in termini di frequentazione, ma da una lacuna specifica nella ricerca dedicata all'arte rupestre in quanto tale. Sapere che la frequentazione mesolitica del Cadore era sottostimata non permette di dedurre che lo sia anche l'arte rupestre del Cadore: sono domande separate che richiedono risposte separate. Ma sapere che la ricerca sistematica in una zona tende a trovare quello che cerca, e che la ricerca sistematica specificamente dedicata all'arte rupestre non è mai stata condotta in modo ampio sul territorio veneto, è sufficiente per mantenere aperta la domanda di questo libro, senza presupporne la risposta in nessuna direzione.
